

My results qualitatively aligns with \cite{jacob2003} on how schools with worse quality inflates grade more than good schools. I find that a one unit difference in school quality was associated with a 10 percentage point difference in the extent of inflating grades. \cite{jacob2003} and \cite{denning2023} similarly find concentration of grade manipulation in schools with worse quality. \cite{jacob2003} finds that a one standard deviation in students achievements is associated with a 0.4 percent point probability in manipulating the students scores. \cite{denning2023} finds that a one percentage point increase in shares of students with free lunch was associated with 4.8 more students getting manipulated scores. 

I find evidence on the positive evolution of grade inflation levels consistent with \cite{dee2019} and \cite{bleemer2020gradeinflation}. I find that experiencing a year of grade discretionary grading scheme is associated with on average a 10 percent increase in the average grade inflation level in the following year. This results aligns with\cite{dee2019} which finds that Dartmouths incoming student's average GPA increased by an annual rate of 40 percent, while \cite{bleemer2020gradeinflation} finds the number to be 10 percent. 

My paper provides a novel evidence on the dynamics of grade inflation by decomposition the extent of the grading policy update by their school types and also their inflation levels in the initial year.I find that the private schools schools increased their inflate rate much faster than other type of schools. Private schools increased their inflation rate by 4 folds of the previous year, while difference in inflation levels during the two years were similar for public schools. My findings matches with \cite{bleemer2020gradeinflation} who finds that schools with students with higher income parents exhibit a faster rate of inflation then other types of schools. However, the persistence of grade inflation is confined with subject groups with weaker effect to spreading to less related subjects. \Cref{tab:dynamics} reports that the correlation between same subject groups across the two consecutive years are significantly stronger than the correlation between neighboring subject groups. My findings support how grades are contagious within schools (\cite{https://doi.org/10.1111/j.1468-2354.2007.00454.x}), but also highlights how the spread of the grading policy is varies within a school. 

My results support how students economic affluence influences teachers grading behavior, supporting the findings by \cite{burgess2013} and \cite{d49bd1f0-4cd8-3cd3-b81e-70c074554021}, while contradicting the findings in \cite{hanna2012}. My results point to how teachers, especially at worse quality schools, tended to assign higher grades to students with lower economic affluence, measured by their parents occupational class and the income deprivation index of their permanent residence. \cite{hanna2012} finds that lower caste students were assigned at most 0.9 standard deviation lower grades than a student from a higher caste. 

I provide evidence that grade inflation is positively correlated with the academic achievements of the student, measured by their GCSE scores. However, the monotone pattern is observed mainly in low quality schools. The extent of the grade inflation is exponentially larger for students with higher GCSE math scores. The results matches \cite{cornwell2013} and \cite{d49bd1f0-4cd8-3cd3-b81e-70c074554021}, who finds similar effect in primary school children. \cite{cornwell2013} find that an one standard deviation increase in math scores increases grade assigned by nearly 0.2
standard deviation points in their grade measure. However, my paper highlights that this pattern is only hold for a subpopulation, namely the low quality schools. For mid-quality schools, the schools/teachers do not seem to assign grades by the academic score of their pupils, and in high quality school grade inflation is larger for the students with lower academic scores. 

I find evidence of gender bias against males. The findings support the findings of the grade bias literature (\cite{lavy2009}, \cite{lavy2018}, \cite{cornwell2013}). Female students received top grades by 6 percentage points higher than male students. The results are at odds with the findings of \cite{hanna2012} and \cite{diamond}, that does not find any biases against males when teachers have grading discretion. \cite{diamond} attempted to identify gender biases by the size of the discontinuity near the cutoff grades for a top grade which is different from my approach with compares final grades between a predicted scores in the standardized regime against the observed teacher-assigned grade. 

I provided a decomposition exercise to explain the changes in the overall grading distribution during the COVID 19 pandemic into the grading biases from the student attributes and the schools grading policies. I find rich heterogeneity in how students attributes lead to reshuffling of top grades within a school with their patterns varying by the schools quality. At low quality schools, students with higher academic achievements increase their chances of receiving a top grade, thereby exacerbating the  achievement inequality across students with varying academic ability. However, low quality schools also provides leniency towards students from socio-economically disadvantaged backgrounds, by assigning better grades towards to the disadvantaged student when comparing academically similar students. However, the mean-spreading property of academic achievements of students weaken in mid-high quality schools, as mid-tier schools show lower tendency to change the ordering of students by academic or demographic characteristics. High quality school inflates the grades for lower academically capable students. However the results for the high quality schools needs to be carefully interpreted as students in high quality schools already have a high chance of receiving a top grade regardless of the grading regime. 

Despite the non-standardized grading regime increased students chances of receiving a top grade, only a limited population of students gained from the opportunity to use the opportunity to secure a placement at a selective university that they had little chance of acceptance. The findings are similar to \cite{dee2019} and \cite{10.1257/pol.20200515} in which they show how improving access to universities don't necessary lead to better placements. My results show that most under-represented students in 2020 did not place into selective universities despite higher grades, even in 2021 although students could have applied to better universities if they successfully predicted that their teachers will grant them higher grades then they deserve. 
